\begin{explanationbox}
\textbf{视角模型：}把 $F_1F_2$ 想成一根固定长度为 $2c$ 的线段，
从椭圆上的点 $P$ 看它，两条视线 $PF_1$、$PF_2$ 的夹角就是
焦角 $\theta$。研究问题不是求某一条焦半径，而是比较这个夹角。

\begin{enumerate}
  \item \textbf{先抓住一个不变量：}
  椭圆定义保证
  \[
  PF_1+PF_2=r_1+r_2=2a.
  \]
  因而焦点三角形 $\triangle PF_1F_2$ 的底边 $2c$、
  周长 $2a+2c$、半周长 $s=a+c$ 都不变。

  \item \textbf{再抓住一个变量：}
  以 $F_1F_2$ 为底边，$P$ 到底边所在直线的高为 $|y_0|$，所以
  \[
  S_{\triangle PF_1F_2}=\frac12\,(2c)|y_0|=c|y_0|.
  \]
  固定周长、固定底边时，三角形面积越大，顶角越大；
  这一关系将在证明中精确化为
  \[
  \tan\frac\theta2=\frac{S}{s(s-2c)}
  =\frac{c|y_0|}{b^2}.
  \]

  \item \textbf{为什么是短轴端点？}
  椭圆方程给出 $0\leq|y_0|\leq b$，
  且只有 $x_0=0$ 时才有 $|y_0|=b$。
  由于 $0\leq\theta/2<90^\circ$，正切函数在这一范围内严格递增，
  因此 $|y_0|$ 越大，焦角就越大。短轴端点是\textbf{距离焦点连线最远}的椭圆点。

  \item \textbf{在短轴端点可以直接看直角三角形：}
  取 $P=(0,b)$，连结原点 $O$，则 $PO\perp F_1F_2$，
  $OF_1=OF_2=c$，$OP=b$，两条焦半径相等。
  $PO$ 平分顶角 $\theta_{\max}$，故
  \[
  \tan\frac{\theta_{\max}}2=\frac{OF_2}{OP}=\frac cb,
  \quad\sin\frac{\theta_{\max}}2=\frac ca=e.
  \]
  \textbf{注意：}这个直角三角形的比例是最大角的\textbf{特殊解法}，
  不能把分母 $b$ 随意换成一般点的 $|y_0|$。

  \item \textbf{题型连接：}
  求焦角范围、判断直角或钝角是否存在、由角反求点的纵坐标、
  通过焦半径与余弦定理计算角度，均可优先使用半角公式。
\end{enumerate}
\end{explanationbox}
